\section{Retos y Casos de Uso}

% --- Ponente 4 ---

\begin{frame}{El Cuello de Botella del Hardware}
\begin{itemize}
    \item Modelos autorregresivos: tardan \textbf{segundos por imagen} en una CPU convencional.
    \item Complejidad computacional $\gg$ decodificadores JPEG.
    \item \textbf{Necesidad absoluta} de hardware dedicado:
    \begin{itemize}
        \item NPU (Neural Processing Unit).
        \item GPU con soporte de inferencia optimizada.
        \item Aceleradores ASIC específicos.
    \end{itemize}
\end{itemize}
\begin{alertblock}{Barrera de adopción}
Sin aceleración hardware, NIC no es viable para decodificación en tiempo real.
\end{alertblock}
\end{frame}

\begin{frame}{Casos de Uso Reales Actuales}
\begin{exampleblock}{Almacenamiento WORM en la nube}
Netflix / Meta: codifican una vez, leen millones de veces. El coste de codificación se amortiza masivamente.
\end{exampleblock}
\begin{itemize}
    \item \textbf{VCM} (Video Coding for Machines):
    \begin{itemize}
        \item Compresión optimizada para visión artificial.
        \item Coches autónomos, vigilancia, drones.
        \item La máquina no necesita ``calidad visual humana''.
    \end{itemize}
    \item Archivos digitales de larga duración.
    \item Streaming adaptativo con decodificación en servidor.
\end{itemize}
\end{frame}

\begin{frame}{Limitaciones Críticas: Alucinaciones Generativas}
\begin{alertblock}{Peligro: contenido inventado}
A tasas de bits extremadamente bajas, los modelos generativos \textbf{inventan texturas realistas pero falsas} (``alucinaciones'').
\end{alertblock}
\begin{itemize}
    \item La imagen reconstruida \textit{parece} correcta, pero contiene detalles ficticios.
    \item \textbf{Inasumible} en dominios críticos:
    \begin{itemize}
        \item Imágenes médicas (diagnóstico erróneo).
        \item Evidencia forense (pruebas judiciales).
        \item Cartografía y teledetección.
    \end{itemize}
    \item Requiere métricas de fidelidad estrictas, no solo perceptuales.
\end{itemize}
\end{frame}
